% Part: normal-modal-logic
% Chapter: tableaux
% Section: simple-S5

\documentclass[../../../include/open-logic-section]{subfiles}

\begin{document}

\olfileid{nml}{tab}{s5}

\olsection{\Log{S5} साठी सोपे \usetoken{P}{tableau}}

\Log{S5} ही प्रणाली वैश्विक प्रतिमानांच्या
वर्गासाठी निर्दोष आणि पूर्ण आहे; अशा प्रतिमानांत
प्रत्येक जग प्रत्येक जगापासून प्राप्य असते.
वैश्विक प्रतिमानांत प्राप्यता संबंध स्वतंत्रपणे
तपासण्याची गरज नसते: ‘$\mSat{M}{!A}[w]$
होईल असे जग~$w$ आहे’ हे विधान तेव्हा आणि
तेव्हाच खरे असते जेव्हा असे~$w$ हे~$u$
पासून प्राप्य असते. म्हणून \Log{S5} मध्ये
प्रतिमानाची व्याख्या केवळ जगांचा संच आणि
सत्य-मूल्यांकन~$V$ यांद्वारे करता येते. यावरून
!!{tableau} चे नियमही सोपे करता येतील असे
सुचते. सर्वसाधारण प्रकरणात पूर्वचिन्हे म्हणून
धन पूर्णांकांच्या क्रमिका वापरतो, म्हणजे कोणत्या
पूर्वचिन्हाने नाव दिलेले जग दुसऱ्या जगापासून
प्राप्य आहे याची नोंद ठेवता येते:
$\sigma.n$ हे~$\sigma$ पासून प्राप्य
जगाचे नाव देते. पण \Log{S5} मध्ये कोणतेही
जग कोणत्याही जगापासून प्राप्य असते, त्यामुळे
अशी नोंद ठेवण्याची गरज नाही. त्याऐवजी
धन पूर्णांक पूर्वचिन्हे म्हणून वापरता येतात.
सोपे केलेले नियम \olref{tab:rules-S5}
मध्ये दिले आहेत.

\begin{table}
  \begin{center}
    \def\arraystretch{3}\def\fCenter{}
    \begin{tabular}{|c|c|}
    \hline
    \iftag{prvBox}{
      \AxiomC{\sFmla{\True}{\Box !A}[n]}
      \RightLabel{\TRule{\True}{\Box}}
      \UnaryInfC{\sFmla{\True}{!A}[m]}
      \DisplayProof
      &
      \AxiomC{\sFmla{\False}{\Box !A}[n]}
      \RightLabel{\TRule{\False}{\Box}}
      \UnaryInfC{\sFmla{\False}{!A}[m]}
      \DisplayProof\\
      $m$ वापरलेले आहे & $m$ नवे आहे
      \\[1ex]
      \hline}{}
    \iftag{prvDiamond}{
      \AxiomC{\sFmla{\True}{\Diamond !A}[n]}
      \RightLabel{\TRule{\True}{\Diamond}}
      \UnaryInfC{\sFmla{\True}{!A}[m]}
      \DisplayProof
      &
      \AxiomC{\sFmla{\False}{\Diamond !A}[n]}
      \RightLabel{\TRule{\False}{\Diamond}}
      \UnaryInfC{\sFmla{\False}{!A}[m]}
      \DisplayProof\\
      $m$ नवे आहे & $m$ वापरलेले आहे
      \\[1ex]
      \hline}{}
    \end{tabular}
  \end{center}
  \caption{\Log{S5} साठी सोपे केलेले नियम.}
  \ollabel{tab:rules-S5}
\end{table}

\begin{ex}
  $\Log{S5} \Proves \Ax{5}$, म्हणजे
  $\Diamond!A \lif \Box\Diamond!A$, हे दाखवणारा
  सोपा केलेला बंद टॅब्लो पुढीलप्रमाणे आहे.
  \begin{oltableau}
    [\pFmla{\False}{\Diamond\formula{A} \lif \Box\Diamond \formula{A}}{1},
      just = \TAss
      [\pFmla{\True}{\Diamond \formula{A}}{1}, just = {\TRule{\False}{\lif}[1]}
        [\pFmla{\False}{\Box\Diamond \formula{A}}{1},
          just = {\TRule{\False}{\lif}[1]}
          [\pFmla{\False}{\Diamond \formula{A}}{2},
            just = {\TRule{\False}{\Box}[3]}
            [\pFmla{\True}{\formula{A}}{3},
              just = {\TRule{\True}{\Diamond}[2]}
                [\pFmla{\False}{\formula{A}}{3},
                  just = {\TRule{\False}{\Diamond}[4]}, close]
            ]
          ]
        ]
      ]
    ]
  \end{oltableau}
\end{ex}

\end{document}
